%
% $Id: $
%
%
% Compilar a .pdf con LaTeX (pdflatex)
% Es necesario instalar Beamer (paquete latex-beamer en Debian)
%

%
% Gr�ficos:
% Los gr�ficos pueden suministrarse en PNG, JPG, TIF, PDF, MPS
% Los EPS deben convertirse a PDF (usar epstopdf)
%

\documentclass{beamer}
\usetheme{Warsaw}
\usebackgroundtemplate{\includegraphics[width=\paperwidth]{format/libresoft-bg.png}}
\usepackage[spanish]{babel}
\usepackage[latin1]{inputenc}
\usepackage{graphics}
\usepackage{amssymb} % Simbolos matematicos

%\definecolor{libresoftgreen}{RGB}{162,190,43}
%\definecolor{libresoftblue}{RGB}{0,98,143}

%\setbeamercolor{titlelike}{bg=libresoftgreen}

%% Metadatos del PDF.
\hypersetup{
  pdftitle={About GSyC/LibreSoft},
  pdfauthor={Jesus M. Gonzalez-Barahona, Gregorio Robles},
  pdfcreator={GSyC/LibreSoft, Universidad Rey Juan Carlos},
  pdfproducer=PDFLaTeX,
  pdfsubject={Systems Course},
}
%%

\begin{document}

\title{About GSyC/LibreSoft}
\subtitle{Presentation of the Systems Course}
\author{Gregorio Robles}
\institute{grex@gsyc.urjc.es \\
GSyC/LibreSoft, Universidad Rey Juan Carlos}

\date{Madrid, March 25th 2011}

\frame{
\maketitle
\begin{center}
\includegraphics[width=6cm]{format/gsyc-urjc}
\end{center}
}


% Si el titulo o el autor se quieren acortar para los pies de p�gina
% se pueden redefinir aqu�:
%\title{Titulo corto}
%\author{Autores abreviado}


% LICENCIA DE REDISTRIBUCION DE LAS TRANSPAS
\frame{
~
\vspace{4cm}

\begin{flushright}
\copyright 2002-2011 Jesus M. Gonzalez-Barahona, Gregorio Robles. \\

Some rights reserved. This presentation is distributed under the
``Attribution-ShareAlike 3.0'' license, by Creative Commons,
available at \url{http://creativecommons.org/licenses/by-sa/3.0/}

\end{flushright}
}


%\section{}


%%---------------------------------------------------------------

\begin{frame}
\frametitle{The Universidad Rey Juan Carlos}

\begin{itemize}
\item Young public university
\item Four \emph{campi} in Madrid: Vic�lvaro, Alcorc�n, M�stoles and Fuenlabrada
\item Around 2,000 staff (1,200 academics); 27,000 students
\item Structured in faculties: ours is the \emph{Escuela Superior de Ingenieros
de Telecomunicaci�n}
\item Official degrees: now moving to the Bologna system
\begin{itemize}
  \item Degree (240 ECTS)
  \item Master (60-120 ECTS)
  \item PhD
\end{itemize}
\item Own titles
\end{itemize}

\begin{flushright}
\url{http://www.urjc.es/}
\end{flushright}

\end{frame}


%%---------------------------------------------------------------

\begin{frame}
\frametitle{Madrid On Rails (MoR)}

\begin{itemize}
\item Main goal: disseminate and promote the use of emerging technologies within small and medium companies in Madrid.
\item Specific goals:
\begin{itemize}
  \item Become a center of international reference in technology, emerging applications and services based on Free Software.
  \item Undertake innovation projects of national and international collaboration with universities, research institutes and companies
  \item Provide specialized training to SMEs in the emerging disciplines of free software and web 2.0, both onsite and online.
\end{itemize}
\end{itemize}

\begin{flushright}
\url{http://madridonrails.com/}
\end{flushright}

\end{frame}



%%---------------------------------------------------------------

\begin{frame}
\frametitle{The GSyC/LibreSoft group}

\begin{itemize}
\item Research group in Uni. Rey Juan Carlos (Fuenlabrada, Spain)
\item About 25 researchers and support staff (including students)
\item Focus on libre (free, open source) technologies, from several
  points of view
\item Backgrounds on software development, data communications,
  telecommunications, etc.
\item Involvement in 6th and 7th FP projects (CALIBRE, FLOSSInclude,
  FLOSSMETRICS, QUALOSS, QUALIPSO, Alert), national projects (Morfeo,
  Vulcano), projects
  funded by companies (Vodafone, Telefonica, Infojobs, Google, Andago) and competitive
  tenders (European Commission, CENATIC)
\end{itemize}

\begin{flushright}
\url{http://libresoft.es/}
\end{flushright}

\end{frame}



%%---------------------------------------------------------------

\begin{frame}
\frametitle{Main activity lines (GSyC/LibreSoft)}

\begin{itemize}
\item Study of libre software development (quantitative, empirical
  approach)
\item Study of Wikipedia and other on-line communities
\item Study and assessment of libre software communities (specially
  those promoted by companies, or composed in part by companies)
\item Technologies for software development (specially those designed
  for practices found in libre software projects)
\item Integration of mobile devices with the Internet
\item Augmented reality in several fields
\item Web 2.0 technologies and the social Internet
\item {\bf Server architecture with free software: Virtualization, clustering and cloud computing}
\end{itemize}

\end{frame}


%%---------------------------------------------------------------

\begin{frame}
\frametitle{Once upon a time...}

\begin{itemize}
\item Interest in libre software (since early 1990s) \\
  e.g: PDSoft Interest Group in Madrid (circa 1992)
\item Collaboration in the community (since mid 1990s) \\
  e.g.: Talks a-la-carte, participation in projects,
\item Education and training on libre software and libre software
  technologies (since mid 1990s) \\
  e.g.: Use of NetBSD in courses (circa 1993)
\item Research on libre software (since early 2000s) \\
  e.g.: Counting potatoes (circa 2001)
\end{itemize}

\end{frame}

%%---------------------------------------------------------------

\begin{frame}
\frametitle{Towards this course}

\begin{itemize}
\item GSyC/LibreSoft research group\\
  (since 2003)
\item Libre software PhD subject\\
  (since 2003)
\item ``Introducci�n al software libre''\\
  (since 2004)
\item MoLOS group, design of an European master\\
  (since 2004)
\item Master on libre software, thanks to Caixa Nova, Igalia
  \\(since 2007)
\item Master edition in Madrid\\
  (since 2008)
\item Bologna-master (``m�ster universitario'') in Madrid\\
  (since 2010)
\item Android and Systems course in Madrid\\
  (since 2011)
\end{itemize}

\end{frame}

%%---------------------------------------------------------------

\begin{frame}
\frametitle{Education on libre software}

\begin{itemize}
\item Libre software has gained a lot of attention from companies,
  public administrations, users of software in general
\item Libre software is a complex (and new) world
\item Many experts on libre software are needed
\item Systems have always been related to libre software
\end{itemize}

\end{frame}

%%---------------------------------------------------------------

%\begin{frame}
%\frametitle{Other participants and supporters}

%\begin{itemize}
%\item Caixa Nova: made this possible
%\item Igalia: good real world experience about libre software
%\item CENATIC: promotion of libre software practices
%\item Medialab Prado: new approaches to libre software
%\end{itemize}

%Prospective supporters: ATI, Morfeo Community. 
%Some other supporters being confirmed.

%\begin{flushright}
%\url{http://master.libresoft.es/}
%\end{flushright}

%\end{frame}

%%---------------------------------------------------------------

\begin{frame}

\begin{center}
{\Huge Come and enjoy!}
\end{center}

\end{frame}


\end{document}
